% !TEX program = lualatex
\documentclass{standalone}
\usepackage{pgfplots}
\usepgfplotslibrary{patchplots,fillbetween}
\pgfplotsset{compat=1.18}
\usepackage[dvipsnames]{xcolor}
\usepgfplotslibrary{colormaps}
\newif\ifexperimental
\newif\iferror

\experimentaltrue % Set to true to include experimental curves
\errorfalse % Set to true to include error bars

\pgfplotsset{
/pgfplots/colormap={blues}{rgb255=(66,106,188) rgb255=(255,255,255)}
}

\pgfplotsset{
    contour/every contour label/.style={
        sloped,
        transform shape,
        inner sep=2pt,
        every node/.style={black,fill=none},
        /pgf/number format/relative={\pgfplotspointmetarangeexponent},
    },
} 

% \def\plotcountourline#1#2{
% 	\addplot3[
% 		point meta=explicit,
% 		contour lua={levels={#1},
% 				labels=true,
% 				draw color=\colorE,
% 				labels over line,
% 				label distance=#2, % change this to ajust numbers
% 				every contour label/.append style={
% 						every node/.append style={\colorE},
% 					},				},
% 		thick,
% 		restrict x to domain=0:70,
% 		restrict y to domain=0:4,
% 	]
% 	table [ header=false,   x index=0,
% 			y index=1,
% 			meta index=2] {data/DeltaN.dat};
% }

\def\plotcontourline#1#2#3{%
    % #1 = contour level
    % #2 = x-range for placing label (min:max)
    % #3 = y-range for placing label (min:max)
    %
    % First: draw full contour line (no labels)

    % Second: overlay same contour but restricted in x/y to force label at desired place
    \addplot3[
        contour lua={levels={#1},
            labels over line,
            label distance=50pt,
            every contour label/.append style={every node/.append style={\colorE}},
        },
        restrict x to domain=#2,
	restrict y to domain=#3,
    ]
    table [header=false, x index=0, y index=1, meta index=2] {data/DeltaN.dat};
    \addplot3[
    contour lua={levels={#1}, labels=false, draw color=\colorE},
    thick,
    restrict x to domain=0:70,
    restrict y to domain=0:4,
    ]
    table [header=false, x index=0, y index=1, meta index=2] {data/DeltaN.dat};
}

\begin{document}
\begin{tikzpicture}[scale=1.0, every node/.style={scale=1.2}]
	\begin{axis}[
			xlabel={$w/\delta^*$},
			ylabel={$d/\delta^*$},
			% small,
			% view={0}{90},
			colorbar,
			colorbar style={
					title={$\Delta n$},     % <-- Your label here
					title style={at={(1.8,0.5)},anchor=west}, % Shift label down if needed
				},
			colormap name=blues,
			ytick distance=0.5,
			ylabel style={rotate=-90, at={(-0.08,0.5)}},
			% unbounded coords=jump,  % This hides values > 5000
			grid=both,
			restrict x to domain*=0:70,
			restrict y to domain*=0:4,
			% point meta min=0.5,
			point meta max=6,
			legend style={
					legend cell align=left,
				},
			scale only axis,
			axis on top,
			set layers,
			enlargelimits=false,
			view={0}{90},
		]

		\def\xmax{\pgfkeysvalueof{/pgfplots/xmax}}
		\def\xmin{\pgfkeysvalueof{/pgfplots/xmin}}
		\def\colorA{Cyan}
		\def\colorB{YellowOrange}
		\def\colorC{gray!80!black}
		\def\colorD{Magenta}
		\def\colorE{Blue}
		\def\size{2pt}
		\def\thickness{very thin}
		\def\thicknessline{very thick}
		\def\linestyle{densely dashed}
		\def\heightRef{0}
		\def\mark{*}
		% standar layers in order are:
		% 	axis background,axis grid,axis ticks,axis lines,axis tick labels, pre main,main,axis descriptions,axis foreground

		%%%%%%%%%%%%%%%%%%%%%%%%%%%%%%%%%%%
		% Filled contour plot
		\pgfonlayer{axis background}
		\addplot[point meta=explicit,
			surf,
			shader=flat,
			draw=none,
			% contour filled={
			% 		% numbers=14,
			% 		levels={0.5,1,...,6},
			% 	},
			contour/labels=true, % enable label drawing
		]
		table [ header=false,   x index=0,
				y index=1,
				meta index=2] {data/DeltaN.dat};
		% \addplot3[
		% 	point meta=explicit,
		% 	contour lua={levels={1,2,3,4},
		% 			labels=true,
		% 			draw color=\colorE,
		% 			labels over line,
		% 			label distance=300pt, % change this to ajust numbers
		% 			every contour label/.append style={
		% 					every node/.append style={\colorE},
		% 				},				},
		% 	thick,
		% 	restrict x to domain=0:70,
		% 	restrict y to domain=0:4,
		% ]
		% table [ header=false,   x index=0,
		% 		y index=1,
		% 		meta index=2] {data/DeltaN.dat};
		% \addplot3[
		% 	point meta=explicit,
		% 	contour lua={levels={5},
		% 			labels=true,
		% 			draw color=\colorE,
		% 			labels over line,
		% 			label distance=48pt, % change this to ajust numbers
		% 			every contour label/.append style={
		% 					every node/.append style={\colorE},
		% 				},				},
		% 	thick,
		% 	restrict x to domain=0:70,
		% 	restrict y to domain=0:4,
		% ]
		% table [ header=false,   x index=0,
		% 		y index=1,
		% 		meta index=2] {data/DeltaN.dat};
		% \plotcountourline{1}{250pt};
		% \plotcountourline{2}{200pt};
		% \plotcountourline{3}{150pt};
		% \plotcountourline{4}{150pt};
		% \plotcountourline{5}{10pt};
		\plotcontourline{1}{10:20}{0.5:1};	
		\plotcontourline{2}{24:26}{0.5:1};
		\plotcontourline{3}{35:45}{0.75:1.25};
		\plotcontourline{4}{53:55}{0.75:1.5};
		\plotcontourline{5}{68:70}{1:1.5};
		
		\endpgfonlayer
		\addplot[
			smooth,
			draw=none,
			name path global=bifurcation,
		] coordinates {
				(16.25, 4.0)
				(15.75, 3.75)
				(15.25, 3.5)
				(15.25, 3.25)
				(15.25, 3.0)
				(15.25, 2.75)
				(20.5, 2.5)
				(31.5, 2.25)
				(39.0, 2.0)
				(46.0, 1.75)
				(56.5, 1.5)
				(65.0, 1.375)
				(70.0, 1.375)
				(75.0, 1.375)
				(80.0, 1.375)
				(100.0, 1.375)
				(120.0, 1.375)
				(130.0, 1.375)
			};
		% 2. Define a horizontal path at ymax
		\addplot[
			name path=top,
			draw=none
		] coordinates {
				(16.25, 4)
				(15.75, 4)
				(15.25, 4)
				(15.25, 4)
				(15.25, 4)
				(15.25, 4)
				(20.5, 4)
				(31.5, 4)
				(39.0, 4)
				(46.0, 4)
				(56.5, 4)
				(65.0, 4)
				(70.0, 4)
				(75.0, 4)
				(80.0, 4)
				(100.0, 4)
				(120.0, 4)
				(130.0, 4)
			};

		\addplot[
			white,
		] fill between[
				of=top and bifurcation,
			];

		%%%%%%%%%%%%%%%%%%%%%%%%%%%%%%%%%%%%%%%

		%%%%%%%%%%%%%%%%%%%%%%%%%%%%%%%%%%%
		% Bifurcation path 
		\addplot[
			smooth,
			no marks,
			draw=\colorC!70!black,
			\thicknessline,
			\linestyle,
		] coordinates {
				(16.25, 4.0)
				(15.75, 3.75)
				(15.25, 3.5)
				(15.25, 3.25)
				(15.25, 3.0)
				(15.25, 2.75)
				(20.5, 2.5)
				(31.5, 2.25)
				(39.0, 2.0)
				(46.0, 1.75)
				(56.5, 1.5)
				(59.0, 1.46)
				(65.0, 1.375)
				(70.0, 1.375)
				(75.0, 1.375)
				(80.0, 1.375)
				(100.0, 1.375)
				(120.0, 1.375)
				(130.0, 1.375)
			};
		%%%%%%%%%%%%%%%%%%%%%%%%%%%%%%%%%%%%%%%

		% level sets experimental curves
		% I fitted the curves by nake eye using geogebra
		\ifexperimental
			\pgfonlayer{axis background}
			\addplot3 [
				contour prepared={labels=true,
						draw color=\colorD,
						labels over line,
						label distance=250pt, % change this to ajust numbers
						every contour label/.append style={
								every node/.append style={\colorD},
							},				},
				thick,
				restrict x to domain=0:70,
				restrict y to domain=0:4,
				legend image post style={draw=red, thick},
			] table [header=false,   x index=0,
					y index=1,
					meta index=2]{data/expData.csv};
			\endpgfonlayer

			\iferror
				\addplot[no marks,
					draw=\colorD!81!black,
					\thicknessline,
					error bars/.cd,
					x dir=both, x explicit,
					error bar style={\colorD!81!black, thick},
				] table [
						x index=0,
						y index=1,
						x error minus index=2,
						x error plus index=3,
						header=false,
					] {data/expDataErr_1.csv};

				\addplot[no marks,
					draw=\colorD!81!black,
					\thicknessline,
					error bars/.cd,
					x dir=both, x explicit,
					error bar style={\colorD!81!black, thick},
				] table [
						x index=0,
						y index=1,
						x error minus index=2,
						x error plus index=3,
						header=false,
					] {data/expDataErr_2.csv};
				\addplot[no marks,
					draw=\colorD!81!black,
					\thicknessline,
					error bars/.cd,
					x dir=both, x explicit,
					error bar style={\colorD!81!black, thick},
				] table [
						x index=0,
						y index=1,
						x error minus index=2,
						x error plus index=3,
						header=false,
					] {data/expDataErr_3.csv};
				\addplot[no marks,
					draw=\colorD!81!black,
					\thicknessline,
					error bars/.cd,
					x dir=both, x explicit,
					error bar style={\colorD!81!black, thick},
				] table [
						x index=0,
						y index=1,
						x error minus index=2,
						x error plus index=3,
						header=false,
					] {data/expDataErr_4.csv};
			\fi

			% \addplot[no marks,
			% smooth,
			% draw=\colorD!80!black,
			% \thicknessline,
			% \linestyle,
			% ] table [
			% x index=0,
			% y index=1,
			% header=false,
			% ] {data/expData_5.csv} node[pos=0.5,anchor=west,\colorD!80!black]{$5$};
			% \addplot[no marks,
			% smooth,
			% draw=\colorD!80!black,
			% \thicknessline,
			% \linestyle,
			% ] table [
			% x index=0,
			% y index=1,
			% header=false,
			% ] {data/expData_6.csv} node[pos=0.5,anchor=west,\colorD!80!black]{$6$};
		\fi

		\legend{,,,,,,,,,,,,,,
			{Transition boundary},{Crouch et al. (2022)},
		}

	\end{axis}
\end{tikzpicture}
\end{document}
